\subsection*{Design and evaluation of the phenomic-diversity recovery simulations}
\phantomsection
\label{sec:neff-supplement}

This section describes the experiment summarized in
\hyperref[sec:neff-main]{Simulation-based recovery of phenomic diversity}.
The simulation design, seeds, complete per-trial outputs, and figure scripts
are provided in \texttt{validation/}\allowbreak\texttt{neff-recovery}.

\paragraph{Tree-generating process and conditioning.}
We simulated 100 independent, fully sampled, bifurcating pure-birth trees,
each starting with two lineages at a crown age of 300~Ma. Speciation rates
were $0.010$, $0.028$, and $0.010~\mathrm{Myr}^{-1}$ during the intervals
300--180, 180--80, and 80--0~Ma, respectively; extinction was zero throughout.
The elevated rate applied to every lineage in the middle interval.
Thus the model is a time-dependent, state-independent special case of BiSSE,
with equal speciation rates for the two states within each epoch. The binary
low/high regime indicator was determined by time, not simulated as an
independently evolving character. We generated trees by exact exponential
waiting times within epochs. \texttt{diversitree::make.bisse.td} was used to
verify a finite likelihood under the generating parameters, not to generate
trees or estimate diversification parameters.

Tree proposals were retained if they contained 120--250 tips, all branch
lengths exceeded $10^{-5}$~Myr, and the focal-branch design could be realized.
Candidate temporal anchors were spaced 1~Myr apart within the elevated-rate
interval, leaving a 3-Myr margin at either end. A valid anchor intersected
at least 20 lineages and had at least ten eligible branches spanning the
6-Myr interval centred on the anchor. Eligible branches subtended at least
six tips, left at least five tips outside their descendant clade, and spent
at least half their duration in the elevated-rate interval. We selected a
valid anchor uniformly, then selected ten eligible branches without
replacement. Proposals with more than 250 tips were discarded in full;
other proposals failing the stated conditions were also discarded, with a
maximum of 1,000 proposals per trial. These conditions depend only on the
tree and focal-branch geometry. No acceptance criterion used character
states, reconstructed rates, $N_{\mathrm{eff}}^{+}$, or correlation results.
The experiment is consequently conditional on these size and geometry criteria.

\paragraph{Character origins and subsequent evolution.}
Each focal branch defined one anatomical unit with 20 separately generated
binary characters, giving 200 characters in ten units per tree. For each
character, the origin was drawn uniformly within the portion of its focal
branch lying between 180 and 80~Ma. At that instant, the character changed
from 0 to 1. Evolution then continued along the remaining focal branch and
all descendant branches under an equal-rates continuous-time Markov model,
$q_{01}=q_{10}=10^{-4}~\mathrm{Myr}^{-1}$, including after the elevated
speciation interval ended. States before the origin and outside its
descendant clade remained~0. We retained exact event times, branch identities,
and simulated histories, including the forced origin and all subsequent
losses or regains. Each character was generated once, without rejection
based on its terminal state frequencies. Conditional on the tree and origin,
the expected number of subsequent transitions is the common transition rate
multiplied by the total descendant branch duration exposed after origin.

\paragraph{Reconstruction and constant observations.}
Reconstruction received the fixed tree and terminal observations only.
It did not receive focal-branch identities, origin times, generating transition
rates, or simulated internal histories. Each variable character was fitted
with an equal-rates Mk model using \texttt{ontophylo} development version~0.3.0,
with three optimization restarts and the \texttt{maddfitz} root treatment
\citep{fitzjohn2009}. The package fits models through \texttt{corHMM} and
passes the fitted transition matrix and root treatment to \texttt{phytools}
for conditional stochastic mapping \citep{corHMM,revell2024}.
We generated 100 maps per variable character, holding its fitted model and
transition matrix fixed. Map intervals therefore describe conditional
history uncertainty and exclude parameter, model-selection, and tree uncertainty.

Some characters can lose state~1 from all surviving tips despite having
undergone a true origin. As in the taxon-sampling experiment, constant
observations were recorded and omitted from reconstruction because the
available observed-state fitting implementation cannot estimate their
binary model. They were not replaced by zero-change maps or redrawn.
All ten units were retained; character retention was reported per unit and
trial. Primary recovery comparisons used the complete simulated histories
as truth. A second comparison used only the simulated histories of characters
retained for reconstruction, separating the effect of lost terminal
variation from reconstruction error. The batch contained
34 constant characters in 28
trials; 100 trials yielded reconstructed trajectories.
All prespecified trials remain in the reported denominators.

\begin{samepage}
\paragraph{Rates, acceleration patterns, and effective diversity.}
For anatomical unit $u$, branch $b$, and mapping replicate $m$, we summed
mapped transitions across its retained characters and divided by branch
duration $\ell_b$:
\[
  r_{bum}=\frac{\sum_{c\in u}n_{bcm}}{\ell_b}.
\]
\end{samepage}
We did not divide by character number. These are realized branch change
rates, not independently fitted branch-specific transition matrices.
For each unit and mapping replicate separately, a branch was classified as
accelerated if its rate was strictly greater than the 95th percentile of
that unit's rates across all branches, using R's type-7 quantile.
Constant rate vectors were left unclassified. This strict rule avoids
classifying zero-change branches as accelerated when the percentile is zero.
We also calculated the original inclusive-cutoff variant as a secondary
sensitivity analysis, retaining undefined correlations when its trajectory
was constant. In this batch, all 100 true-history and all 100 reconstructed
inclusive-cutoff trajectories were constant, giving no estimable
correlations for that variant.

At each 1-Myr midpoint from 299.5 to 0.5~Ma, we selected all extant branches
and represented each by its joint binary acceleration pattern across the
ten units. If $f_{tm}$ is the fraction with at least one accelerated unit and
$p_{ktm}$ is the frequency of nonzero pattern $k$ among these branches,
\[
  N_{\mathrm{eff},tm}^{+}
  = f_{tm}\exp\left(-\sum_k p_{ktm}\log p_{ktm}\right).
\]
The metric was zero when no branch was accelerated. We calculated it within
each mapping replicate before averaging across maps; thresholding averaged
rates or treating marginal unit frequencies as independent patterns would
not reproduce this calculation. Applying the same operations to the actual
simulated histories provided the true-history reference.
Because rates summarize whole branches, even the true-history rate metric
can extend outside the precise interval containing individual origin events.

Marginal character rates were means across branches present at each time
slice. Marginal diversification was the corresponding mean of the known
generating $\lambda-\mu$; with equal rates among contemporaneous lineages,
it equals the specified epoch rate. The simulation therefore tests recovery
of the phenomic component against known diversification, rather than joint
recovery of two estimated trajectories. Figure~\ref{fig:neff-example} shows
these distinct sources for prespecified trial~001, and
Figure~\ref{fig:neff-unit} links unit-level branch classifications to the
reconstructed diversity trajectory.

\paragraph{Recovery and correlation summaries.}
We quantified true--reconstructed agreement by the trajectory correlation,
root-mean-square error, mean bias, and peak age on the 300-point time grid.
These are descriptive recovery measures. For association tests, we averaged
$N_{\mathrm{eff}}^{+}$ and diversification within the same 15 complete,
non-overlapping 20-Myr bins and used a two-sided Pearson correlation test.
We tabulated the numbers and percentages with $p<0.05$, $p<0.01$, and
$p<0.001$, both irrespective of sign and requiring a positive correlation.
Percentages use all 100 prespecified trials; the accompanying tables also
provide denominators restricted to estimable correlations and explicitly
list unavailable results. Exact 95\% binomial intervals quantify Monte Carlo
uncertainty in these frequencies. They do not correct the individual
correlation tests for dependence among time bins.

Across-trial trajectory bands show the central 95\% range of trial means,
not confidence intervals for the grand mean. Within-trial bands show the
central 95\% range across conditional stochastic maps. This experiment
contains an imposed-association scenario only: it does not estimate null
false-positive rates, address extinct lineages or incomplete sampling, or
establish a causal effect of character evolution on diversification.

\paragraph{Computation and verification.}
Calculations ran on CSC Roihu with fixed trial-specific seeds and a frozen
\texttt{ontophylo} source snapshot. Each tree and dataset was independent;
trial~001 served as a technical preflight and remained part of the 100-trial
batch. The gate checked computation and metric correctness, not the size or
significance of correlations. Character fits and mapping batches were
checkpointed. Independent checks reconstructed the strict acceleration
patterns and effective-diversity formula directly from branch-rate arrays,
verified true rates against simulated event counts, and recalculated the
reported Pearson tests. Package versions, source checksums, fitting and
mapping statuses, and cluster accounting are retained with the analysis.
